\subsection{艾森斯坦多项式}

定义 2. --- 设 A 是环，p 是 A 的素理想，P 是 A{[}T{]} 中的多项式。如果 P 满足以下条件，则称 P 是关于 p 的艾森斯坦多项式：

\begin{enumerate}
\def\labelenumi{\alph{enumi})}
\item
  P 是次数 \(d \geqslant 1\) 的首一多项式；
\item
  有 \(\mathrm{P(T)}\equiv \mathrm{T}^d\bmod \mathfrak{pA}[\mathrm{T}];\)
\item
  有 \(\mathbf{P}(0)\notin \mathfrak{p}^2\)
\end{enumerate}

换言之，关于 p 的艾森斯坦多项式是形如 \(\mathbf{P}(\mathbf{T}) = \mathbf{T}^{d} + \sum_{i=1}^{d} a_{i} \mathbf{T}^{d-i}\) 的多项式，其中 \(d \geqslant 1\)，\(a_{1}, \ldots, a_{d-1}\) 属于 p，且 \(a_{d}\) 属于 \(p - p^{2}\)。

称 P 是关于 \(pA_{p}\) 的艾森斯坦多项式，当且仅当 P 在多项式环 \(A_{p}[T]\) 中的典范像是关于理想 \(pA_{p}\) 的艾森斯坦多项式；这也意味着 P 是关于 p 的艾森斯坦多项式，并且还满足下述比 c) 更强的条件：

\(c^{\prime})\) A 中每个元素 \(a\)，只要 \(a\mathbf{P}(0)\in\mathfrak{p}^{2}\)，就属于 \(\mathfrak{p}\)。

命题 3. --- 设 A 是环，p 是 A 的素理想，P ∈ A{[}T{]} 是关于 p 的艾森斯坦多项式。

\begin{enumerate}
\def\labelenumi{\alph{enumi})}
\item
  不存在形如 \(P = P_{1}P_{2}\) 的分解，其中 \(P_{1}\) 和 \(P_{2}\) 是 A{[}T{]} 中不同于 1 的两个首一多项式。
\item
  假设 A 整闭，分式域为 K。则 P 在 K[T] 中不可约。
\end{enumerate}

令 \(\varphi\) 为从 A 到 A/p 的分式域 \(k\) 的典范同态，并令 \(\varphi':\mathrm{A}[\mathrm{T}]\to k[\mathrm{T}]\) 为 \(\varphi\) 的延拓，使得 \(\varphi'(\mathrm{T})=\mathrm{T}\)。假设 \(\mathbf{P}=\mathbf{P}_{1}\mathbf{P}_{2}\)，其中 \(\mathbf{P}_{1}\) 和 \(\mathbf{P}_{2}\) 是 A{[}T{]} 中不同于 1 的两个首一多项式。于是有 \(\mathrm{T}^{d}=\varphi'(\mathbf{P}_{1})\varphi'(\mathbf{P}_{2})\)，且等式成立于 \(k[\mathrm{T}]\) 中，其中 \(d\) 为 P 的次数。若 \(d_{i}\) 是 \(\mathbf{P}_{i}\) 的次数，则有 \(\varphi'(\mathbf{P}_{i})=\mathrm{T}^{d_{i}}\)，即 \(\mathbf{P}_{i}(\mathrm{T})\equiv\mathrm{T}^{d_{i}}\text{ mod. pA[T]}\)，特别地，\(\mathbf{P}_{i}(0)\in\mathfrak{p}\)。但是 \(\mathrm{P}(0)=\mathrm{P}_{1}(0).\mathrm{P}_{2}(0)\) 属于 \(\mathfrak{p}^{2}\)，这与假设矛盾。这证明了 \(a\)。

断言 \(b)\) 由 \(a)\) 以及 V, § 1, no. 3 的命题 11 得出。

设 A 是诺特局部环，且 \(P_{1}, \ldots, P_{r}\) 是 A{[}T{]} 中次数为 \(\geqslant 2\) 的首一多项式。令 q 为 A{[}T \(_{1}\) , \ldots, T \(_{r}\) {]} 的理想，它由 \(\mathrm{P}_{1}(\mathrm{T}_{1}), \ldots, \mathrm{P}_{r}(\mathrm{T}_{r})\) 生成；令 B 为商 A-代数 A{[}T \(_{1}\) , \ldots, T \(_{r}\) {]}/q。对于 \(1 \leqslant i \leqslant r\)，记 \(d_{i}\) 为 \(P_{i}, t_{i}\) 的相关数据：前者表示 P＿i 的次数，后者表示 T \(_{i}\) 模 q 的类；并记 \(\gamma_{i}\) 为 \(c_{i} = \mathrm{P}_{i}(0)\) 模 \(m_{A}^{2}\) 的类。假设有 \(\mathrm{P}_{i}(\mathrm{T}) \equiv \mathrm{T}^{d_{i}} \mod \mathrm{m}_{\mathrm{A}}\mathrm{A}[\mathrm{T}]\)，这对 \(1 \leqslant i \leqslant r\) 成立。

命题 4. --- a) 环 B 是局部诺特环，其极大理想为

\[
\mathfrak {m} _ {\mathrm{B}} = \mathbf {B m} _ {\mathrm{A}} + \sum_ {i = 1} ^ {r} \mathbf {B t} _ {i}.
\]

有 \(\dim(\mathbf{A}) = \dim(\mathbf{B})\)，且 \([\kappa_{\mathbf{B}}: \kappa_{\mathbf{A}}] = 1\)。单项式 \(t_1^{\alpha(1)} \ldots t_r^{\alpha(r)}\)，其中 \(0 \leqslant \alpha(i) < d_i\)，且 \(1 \leqslant i \leqslant r\)，构成 A-模 B 的一个基。

\begin{enumerate}
\def\labelenumi{\alph{enumi})}
\setcounter{enumi}{1}
\item
  令 \(\lambda\) 为从 \(\mathfrak{m}_{\mathbf{A}} / \mathfrak{m}_{\mathbf{A}}^{2}\) 到 \(\mathfrak{m}_{\mathbf{B}} / \mathfrak{m}_{\mathbf{B}}^{2}\) 的典范同态。则 \(\lambda\) 的核是一个 \(\kappa_{\mathbf{A}}\)-向量空间，由 \(\gamma_{1}, \ldots, \gamma_{r}\) 生成。元素 \(t_{1}, \ldots, t_{r}\) 的类关于 \(\kappa_{\mathbf{A}}\) 构成 \(\lambda\) 的余核的一个基。
\item
  B 正则的充要条件是 A 正则，且 \(\gamma_{1},\ldots,\gamma_{r}\) 在 \(\kappa_{\mathbf{A}}\)-向量空间 \(\mathfrak{m}_{\mathbf{A}}/\mathfrak{m}_{\mathbf{A}}^{2}\) 中线性无关。
\end{enumerate}

A-代数 B 同构于张量积 \(B_{1} \otimes_{A} \cdots \otimes_{A} B_{r}\)，其中 \(B_{i} = A[T]/(P_{i})\)，\(1 \leqslant i \leqslant r\)。因此，单项式 \(t_{1}^{\alpha(1)} \ldots t_{r}^{\alpha(r)}\)，其中 \(0 \leqslant \alpha(i) < d_{i}\)，且 \(1 \leqslant i \leqslant r\)，构成 A-模 B 的一个基。特别地，B 在 A 上整，故根据 § 2, no. 3 的定理 1，A 与 B 具有相同维数。

根据 IV, § 2, no. 5 的命题 9 的推论 3，环 B 是诺特的，并且 B 的每个极大理想都包含 B.m\_A。此外，由于对 P\_1, \ldots P\_r 所作的假设以及关系 P\_i(t\_i) = 0，有 t\_i\^{}\{d\_i\} ∈ B.m\_A（1 ≤ i ≤ r）。因此 B 的每个极大理想都包含 t\_1, \ldots, t\_r，因而也包含理想 q' = B.m\_A + Bt\_1 + \ldots{} + Bt\_r。现在有 m\_A = A ∩ q' 且 B = A + q'，故 B/q' 同构于 A/m\_A，而 q' 是 B 的一个极大理想；所以 B 是局部的，且有 {[}κ\_B : κ\_A{]} = 1。这证明了 a)。

置 \(\mathfrak{r} = \mathfrak{m}_{A}^{2} + \sum_{i=1}^{r} A c_i\)，令 \(\varphi\) 为从 \((A/\mathfrak{m}_{A}^{2})[T_1, ..., T_r]\) 到 \(B/\mathfrak{m}_{B}^{2}\) 的典范同态。由于有 \(m_B = B.m_A + \sum_{i=1}^{r} B t_i\)，其核记为 \(n\)；关于 \(\varphi\)，该核由多项式的类 \(\overline{P}_i(T_i)\) 生成，其中多项式为 \(P_i(T_i)\)，并按模 \(\mathfrak{m}_{A}^{2}.A[T_1, ..., T_r]\) 取类；还由单项式 \(T_i T_j\) 和 \(x T_i\) 生成，其中 \(1 \leqslant i, j \leqslant r\)，且 \(x\) 属于 \(m_A/m_A^2\)。根据对 \(P_i\) 所作的假设，即 \(P_i(T) \equiv T^{d_i} \mod. m_A.A[T]\)，可在上述描述中将 \(\overline{P}_i(T_i)\) 替换为 \(\gamma_i\)，而这个描述涉及 \(n\)；因此，环 \(B/m_B^2\) 同构于 \(A/\mathfrak{r}\) \([T_{1}, ..., T_{r}]\) 关于分次理想的商，该理想由单项式 \(T_{i}T_{j}\) 和 \(xT_{i}\) 生成，其中 \(x\) 属于 \(m_A/\mathfrak{r}\)。记 \(\tau_{i}\) 为 \(t_{i}\) 模 \(m_{B}^{2}\) 的类；于是有

\[
\mathrm{B} / \mathfrak {m} _ {\mathrm{B}} ^ {2} = (\mathrm{A} / \mathfrak {r}) \oplus \kappa_ {\mathrm{A}} \tau_ {1} \oplus \dots \oplus \kappa_ {\mathrm{A}} \tau_ {r},\tag{5}
\]

从而

\[
\mathfrak {m} _ {\mathrm{B}} / \mathfrak {m} _ {\mathrm{B}} ^ {2} = (\mathfrak {m} _ {\mathrm{A}} / \mathfrak {r}) \oplus \kappa_ {\mathrm{A}} \tau_ {1} \oplus \dots \oplus \kappa_ {\mathrm{A}} \tau_ {r}.\tag{6}
\]

断言 b) 立即由此得到。

根据公式 (6) 及关系 \([\kappa_{\mathrm{B}}:\kappa_{\mathrm{A}}] = 1\)，有

\[
\left[ \mathfrak {m} _ {\mathrm{B}} / \mathfrak {m} _ {\mathrm{B}} ^ {2}: \kappa_ {\mathrm{B}} \right] = \left[ \mathfrak {m} _ {\mathrm{A}} / \mathfrak {m} _ {\mathrm{A}} ^ {2}: \kappa_ {\mathrm{A}} \right] + \left\{r - \left[ \mathfrak {r} / \mathfrak {m} _ {\mathrm{A}} ^ {2}: \kappa_ {\mathrm{A}} \right] \right\}.\tag{7}
\]

现在，\(\kappa_{\mathrm{A}}\)-向量空间 \(\mathfrak{r}/\mathfrak{m}_{\mathrm{A}}^{2}\) 由 \(\gamma_{1}, \ldots, \gamma_{r}\) 生成，且有

\[
\dim (\mathbf {B}) = \dim (\mathbf {A}) \leqslant [ \mathfrak {m} _ {\mathbf {A}} / \mathfrak {m} _ {\mathbf {A}} ^ {2}: \kappa_ {\mathbf {A}} ].\tag{8}
\]

断言 (c) 立即由公式 (7) 和 (8) 得出。

推论。--- 设 A 是正则诺特局部环，且 P ∈ A{[}T{]} 是关于 m\_A 的艾森斯坦多项式。环 B = A{[}T{]}/(P) 是与 A 同维数的正则诺特局部环，且有 {[}κ\_A : κ\_B{]} = 1。最后，有 m\_B = B.m\_A + Bt，其中 t 是 T 模 (P) 的类。

P 次数至少为 2 的情形由取 r = 1 应用命题 4 得出；当 P 次数为 1，即形如 T - c 且 \(c \in m_{A}\) 时，推论立即得到。

命题 5. --- 设 A 是整环，分式域为 K，L 是 K 的有限代数扩张。记 B 为 L 中 A 的整闭包，p 为 A 的一个素理想。

假设局部环 \(\mathbf{A}_{\mathfrak{p}}\) 是诺特且正则的；设 \(t\) 是 \(\mathbf{L}\) 中的元素，使得 \(\mathbf{L} = \mathbf{K}(t)\)，并假设在 \(\mathbf{A}[\mathbf{T}]\) 中存在元素 \(\mathbf{P}\)，它是关于 \(\mathfrak{p}\mathbf{A}_{\mathfrak{p}}\) 的艾森斯坦多项式，且 \(t\) 是其根。

\begin{enumerate}
\def\labelenumi{\alph{enumi})}
\item
  B 中存在唯一一个位于 p 之上的素理想 q。
\item
  局部环 \(\mathbf{B}_{\mathfrak{q}}\) 是诺特且正则的，并且与 \(\mathbf{A}_{\mathfrak{p}}\) 具有相同维数。
\item
  我们有 \(\mathbf{B}_{\mathfrak{q}} = \mathrm{A}_{\mathfrak{p}}[t]\)。
\item
  从 A/p 到 B/q 的典范同态在这些环的分式域之间诱导出一个同构。
\end{enumerate}

置 \(C = A_{p}[t]\)，记 \(d\) 为 P 的次数。将命题 3 应用于环 \(A_{p}\)，可知艾森斯坦多项式 P 在 K{[}T{]} 中不可约，且 \((1, t, ..., t^{d-1})\) 是 L 关于 K 的一个基，因而也是 C 关于 \(A_{p}\) 的一个基。由于 P 是首一的，从 \(A_{p}[T]\) 到 C 的典范同态的核等于 (P)。根据上面的命题 4 的推论，C 是与 \(A_{p}\) 具有相同维数的正则诺特局部环，C 的极大理想 \(m_{C}\) 由 \(p \cup \{t\}\) 生成，且域 \(\kappa_{C}\) 是 A/p 的分式域的平凡扩张。要证明命题 5，只需证明 B 上存在唯一一个位于 p 之上的素理想 q，并且有 \(C = B_{q}\)。

置 \(S = A - p\)。我们知道（V, § 1, no. 5, 命题 16），其整闭包

\(A_{p}\) 在 L 中等于 \(S^{-1}B\)。此外，t 在 \(A_{p}\) 上整，环 \(C = A_{p}[t]\) 是正则诺特局部环，因而是整闭的（no. 2, 定理 1 的推论 1）。所以有 \(C = S^{-1}B\)。因此，环 \(S^{-1}B\) 是局部的并且只有一个极大理想。根据 V, § 2, no. 1 的命题 1，\(S^{-1}B\) 在 \(pA_{p}\) 之上有唯一的素理想，所以 B 在 p 之上有唯一的素理想 q（同上，引理 1），并且有 \(B_{q} = S^{-1}B = C\)。

推论。--- 假设 \(A_{p}\) 是离散赋值环。则 \(B_{q}\) 是离散赋值环，t 是 \(B_{q}\) 的一致化元，且有

\[
f (\mathrm{B} _ {\mathfrak {q}} / \mathrm{A} _ {\mathfrak {p}}) = 1, e (\mathrm{B} _ {\mathfrak {q}} / \mathrm{A} _ {\mathfrak {p}}) = [ \mathrm{L}: \mathrm{K} ]\tag{9}
\]

(VI, § 8, no. 1).

事实上，离散赋值环正是维数为 1 的正则诺特局部环；置 \(d = [\mathrm{L}:\mathrm{K}]\)，有 \(t^{d}\in\mathfrak{m}_{\mathrm{A}}-\mathfrak{m}_{\mathrm{A}}^{2}\)，从而 \(d=e(\mathrm{B}_{\mathfrak{q}}/\mathrm{A}_{\mathfrak{p}})\)。又有 \([\kappa_{\mathrm{B}}:\kappa_{\mathrm{A}}]=1\)，从而 \(f(\mathrm{B}_{\mathfrak{q}}/\mathrm{A}_{\mathfrak{p}})=1\)。

例子。--- 1) 置 A = Z，L = Q( \(p^{1/d}\) )，其中 p 是素数，d 是满足 d ≥ 2 的整数。记 B 为 L 中 Z 的整闭包。由于多项式 T \(^{d}\) - p 在 Z{[}T{]} 中是关于 pZ \(_{(p)}\) 的艾森斯坦多项式，故 B 中存在唯一一个位于 pZ 之上的素理想 q。因此，存在域 Q( \(p^{1/d}\) ) 的唯一归一化离散赋值 v，使得 v(p) \textgreater{} 0；有 {[}L:K{]} = v(p) = d，且 B/q 是含 p 个元素的域。赋值 v 的环 B\(_{q}\) 等于 Z \(_{(p)}\) {[} \(p^{1/d}\) {]}。

\begin{enumerate}
\def\labelenumi{\arabic{enumi})}
\setcounter{enumi}{1}
\tightlist
\item
  置 A = Z，且 \(L = R_{p^{f}}(\mathbf{Q})\)，其中 p 是素数，f 是满足 \(\geqslant 1\) 的整数（参见 A, V, p. 78）。因此有 \(L = \mathbf{Q}(\zeta)\)，其中 \(\zeta = \exp(2\pi i/p^{f})\)。令 B 为 L 中 Z 的整闭包，P 为 Z{[}T{]} 中满足下式的多项式：
\end{enumerate}

\[
\mathrm{P} (\mathrm{T} - 1) = \left(\mathrm{T} ^ {p^{f}} - 1\right) / \left(\mathrm{T} ^ {p^{f - 1}} - 1\right).
\]

置 \(d = p^{f} - p^{f - 1}\)。有 \(\mathrm{P}(\zeta - 1) = 0\)，\(\mathrm{P}(0) = p\)，且

\[
\mathrm{P} (\mathrm{T} - 1) \equiv (\mathrm{T} - 1) ^ {d} \bmod . p \mathbf {Z} [ \mathrm{T} ],
\]

因而 P(T) ≡ T \(^{d}\)。所以，P 是关于 pZ \(_{(p)}\) 的艾森斯坦多项式。因此，B 中存在唯一一个位于 pZ 之上的素理想 q，并且有 B \(_{q}\) = Z \(_{(p)}\) {[}ζ{]}；此外，ζ - 1 是 B \(_{q}\) 的一致化元，且有

\[
[ \mathrm{L}: \mathrm{K} ] = d = p ^ {f} - p ^ {f - 1}.
\]

若 \(v\) 是 \(\mathbf{Q}(\zeta)\) 的唯一归一化赋值，使得 \(v(p)>0\)，则有 \(v(p)=d\)。此外，域 \(\mathbf{B}/\mathfrak{q}\) 含有 \(p\) 个元素。可以证明（参见 p. 96, exerc. 13），\(\mathbf{B}\) 等于 \(\mathbf{Z}[\zeta]\)。



